\chapter{Dimension Theory}

\section*{HILBERT FUNCTIONS}

Let $A$ be a Noetherian graded ring and $M$ be a finitely-generated graded $A$-module. Let $\lambda$ be an additive function on the class of all finitely generated $A_0$-modules. The \textcolor{orange}{Poincar\'{e} series} of $M$ is the power series
\begin{equation*}
	P(M,t)=\sum_{n=0}^\infty\lambda\left(M_n\right)t^n\qquad\in\mathbb{Z}[[t]].
\end{equation*}

\begin{theorem}{11.1}[\textcolor{red}{Hilbert-Serre theorem}]
	$P(M,t)$ is a rational function in $t$ of the form $f(t)/\prod_{i=1}^d\left(1-t^{k_i}\right)$, where $f(t)\in\mathbb{Z}[t]$.
\end{theorem}

The order of the pole of $P(M,t)$ at $t=1$ is denoted as \textcolor{orange}{$d(M)$}.

\begin{corollary}{11.2}
	If each $k_i=1$, then for all sufficient large $n$, $\lambda\left(M_n\right)$ is a polynomial in $n$. More specifically, if $f(t)=\sum_{k=0}^Na_kt^k$, we have
	\begin{equation*}
		\lambda\left(M_n\right)=\sum_{k=0}^na_k\binom{d+n-k-1}{d-1}.
	\end{equation*}
	This is called \textcolor{orange}{Hilbert function} or \textcolor{orange}{Hilbert polynomial} of $M$, and its leading term is $\left(\sum a_k\right)n^{d-1}/(d-1)!$.
\end{corollary}

\begin{proposition}{11.3}
	If $x\in A_k$ is not a zero-divisor in $M$ then $d(M/xM)=d(M)-1$.
\end{proposition}

\begin{proposition}{11.4}
	Let $A$ be a Noetherian local ring, $\mathfrak{m}$ its maximal ideal, $\mathfrak{q}$ an $\mathfrak{m}$-primary ideal, $M$ a finitely-generated $A$-module, $\left(M_n\right)$ a stable $\mathfrak{q}$-filtration of $M$. Then
	\begin{enumerate}[label=\textup{(\roman*)}]
		\item $M/M_n$ is of finite length for each $n\geqslant0$;
		\item for all suficiently large $n$ this length is a polynomial $\chi_\mathfrak{q}^M(n)$ of degree $\leqslant s$, where $s$ is the least number of generates of $\mathfrak{q}$;
		\item the degree and leading coefficient of $\chi_\mathfrak{q}^M(n)$ depend only on $M$ and $\mathfrak{q}$, not on the filtration chosen.
	\end{enumerate}
\end{proposition}

If $M=A$, the polynomial $\chi_\mathfrak{q}^M(n)$ is called the \textcolor{orange}{characteristic polynomial} of the $\mathfrak{m}$-primary ideal $\mathfrak{q}$, simply noted as $\chi_\mathfrak{q}(n)$.

\begin{proposition}{11.6}
	$\operatorname{deg}\chi_\mathfrak{q}(n)=\operatorname{deg}\chi_\mathfrak{m}(n)$.
\end{proposition}

If $A$ is a ring, $\mathfrak{p}$ a prime ideal in $A$, then the \textcolor{orange}{height} of $\mathfrak{p}$ is defined to be the supremum of chains of prime ideals $\mathfrak{p}_0\subsetneq\mathfrak{p}_1\subsetneq\cdots\subsetneq\mathfrak{p}_r=\mathfrak{p}$ which end at $\mathfrak{p}$. Hence $\operatorname{height}\mathfrak{p}=\dim A_\mathfrak{p}$.

Likewise we define the \textcolor{orange}{depth} of $\mathfrak{p}$, by considering chains of prime ideals which start at $\mathfrak{p}$. Hence, $\operatorname{depth}\mathfrak{p}=\dim A/\mathfrak{p}$.

\section*{DIMENSION THEORY OF NOETHERIAN LOCAL RINGS}

\begin{theorem}{11.14}[Dimension theorem]
	For any Noetherian local ring $A$ the following three integers are equal:
	\begin{enumerate}[label=\textup{(\roman*)}]
		\item the maximum length of chains of prime ideals in $A$ (\textcolor{orange}{Krull dimension});
		\item the degree of the characteristic polynomial $\chi_\mathfrak{m}=l\left(A/\mathfrak{m}^n\right)$;
		\item the least number of generators of an $\mathfrak{m}$-primary ideal of $A$.
	\end{enumerate}
\end{theorem}

\begin{corollary}{11.17}[Krull's principal ideal theorem]
	Let $A$ be a Noetherian ring and let $x$ be an element of $A$ which is nether a zero-divisor nor a unit. Then every minimal prime ideal $\mathfrak{p}$ of $(x)$ has height 1.
\end{corollary}

\begin{corollary}{11.18}
	Let $A$ be a Noetherian local ring, $x$ an element of $\mathfrak{m}$ which is not a zero-divisor. Then $\dim A/(x)=\dim A-1$.
\end{corollary}

\begin{corollary}{11.19}
	Let $\hat{A}$ be the $\mathfrak{m}$-adic completion of $A$. Then $\dim A=\dim\hat{A}$.
\end{corollary}

If $x_1,\cdots,x_d$ generate an $\mathfrak{m}$-primary ideal, and $d=\dim A$, we call $x_1,\cdots,x_d$ a \textcolor{orange}{system of parameters}.

\begin{proposition}{11.20}
	Let $x_1,\cdots,x_d$ be a system of parameters for $A$ and let $\mathfrak{q}=\left(x_1,\cdots,x_d\right)$ be the $\mathfrak{m}$-primary ideal. Let $f\left(t_1,\cdots,t_d\right)$ be a homogeneous polynomial of degree $s$ with coefficients in $A$, and assume that $f\left(x_1,\cdots,x_d\right)\in\mathfrak{q}^{s+1}$. Then all the coefficients of $f$ lie in $\mathfrak{m}$.
\end{proposition}

\section*{REGULAR LOCAL RINGS}

\begin{theorem}{11.22}
	Let $A$ be a Noetherian local ring of dimension $d$, $\mathfrak{m}$ its maximal ideal, $k=A/\mathfrak{m}$. Then the following are equivalent:
	\begin{enumerate}[label=\textup{(\roman*)}]
		\item $G_\mathfrak{m}(A)\cong k[t_1,\cdots,t_d]$ where the $t_i$ are independent indeterminates;
		\item $\dim_k\left(\mathfrak{m}/\mathfrak{m}^2\right)=d$;
		\item $\mathfrak{m}$ can be generated by $d$ elements.
	\end{enumerate}
	The ring satisfying any of the conditions is called \textcolor{orange}{regular} local ring.
\end{theorem}

\begin{proposition}{11.23}
	Let $A$ be a ring, $\mathfrak{a}$ an ideal of $A$ such that $\bigcap_n\mathfrak{a}^n=0$. Suppose that $G_\mathfrak{a}(A)$ is an integral domain. Then $A$ is an integral domain.
\end{proposition}

\begin{proposition}{11.24}
	Let $A$ be a Noetherian local ring. Then $A$ is regular iff $\hat{A}$ is regular.
\end{proposition}

\section*{TRANSCENDENTAL DIMENSION}

Assume that $k$ is an algebraically closed field and let $V$ be an irreducible affine variety over $k$. Let the coordinate ring of $V$ be $A(V)=k[x_1,\cdots,x_n]/I(V)$. Then $k(V)=\operatorname{Frac}A(V)$ is a finitely generated extension of $k$ and so has a finite transcendence degree over $k$, which is defined to be the \textcolor{orange}{dimension} of $V$.

If $P$ is a point with maximal ideal $\mathfrak{m}$ we call the Krull dimension $\dim A(V)_\mathfrak{m}$ the \textcolor{orange}{ local dimension} of $V$ at $P$.

\begin{theorem}{11.25}
	For any irreducible variety $V$ over $k$ the local dimension at any point is equal to its dimension.
\end{theorem}

